% Results-section addendum, placed immediately after Table 2:
\paragraph{Out-of-sample evaluation and ALERTCalifornia leakage.}
The ALERTCalifornia camera network has been deployed and progressively expanded
between 2017 and 2024, with camera siting partially informed by historical
ignition records that include fires from 2021--2023~\cite{alertcalifornia_history,alertcalifornia2026}.
This means that the 2021--2023 columns of Table~\ref{tab:alertcalifornia} probe
the network \emph{partly in-sample} from a siting-decision standpoint. The 2024
column provides the cleanest out-of-sample evaluation, given that the 2024
ignition data were not available when most camera placement decisions were made,
and we therefore use the 2024 dataset for Figure~\ref{fig:alertcalifornia}.

% Discussion paragraph, in the limitations or future-work discussion:
\paragraph{Comparability across baselines.}
A subtlety in comparing our optimised drone network against ALERTCalifornia is
that the two systems were optimised on different information sets. The
drone-network solution is generated entirely from pre-2021 inputs (the static
Pyrologix risk surface and the burnable-area mask), so 2021--2024 outcomes are
fully out-of-sample. The ALERTCalifornia camera placement, by contrast, has been
adjusted over the years partly in response to historical ignitions, and should be
expected to perform best on historical fires. The 2024 column of
Table~\ref{tab:alertcalifornia} is therefore the most directly comparable to the
drone results; a fully controlled comparison would re-optimise the cameras using
only pre-2021 data, which is beyond the scope of this work but constitutes a
natural follow-up.
